\documentclass[11pt,a4paper]{article}
\usepackage[a4paper,hmargin=1in,top=0.75in,bottom=0.75in]{geometry}
\usepackage[T1]{fontenc}
\usepackage{lmodern}
\usepackage{amsmath,amssymb,amsthm,mathtools}
\usepackage{microtype,booktabs,array,enumitem}
\usepackage[svgnames]{xcolor}
\usepackage{framed,thmtools,titlesec,titling,fancyhdr}
\usepackage{etoolbox}
\usepackage[colorlinks=true,linkcolor=blue,citecolor=blue,urlcolor=blue]{hyperref}

% Visual conventions from Furstenberg/frozens; document wording is unchanged.
\renewcommand{\familydefault}{\sfdefault}
\setlength{\parskip}{0.5\baselineskip plus 2pt}
\setlength{\parindent}{0pt}
\setlength{\emergencystretch}{2em}
\titleformat{\section}{\normalfont\sffamily\LARGE\bfseries}{\thesection}{1em}{}
\titleformat{\subsection}{\normalfont\sffamily\Large\bfseries}{\thesubsection}{1em}{}
\titleformat{\subsubsection}{\normalfont\sffamily\Large\bfseries}{\thesubsubsection}{1em}{}
\titleformat{\paragraph}[runin]{\normalfont\sffamily\bfseries}{}{0pt}{}
\pretitle{\begin{center}\sffamily\Huge\bfseries}
\posttitle{\par\end{center}\vskip 0.5em}
\preauthor{\begin{center}\sffamily\LARGE\lineskip 0.5em\begin{tabular}[t]{c}}
\postauthor{\end{tabular}\par\end{center}}
\predate{\begin{center}\sffamily\LARGE}
\postdate{\par\end{center}}
\thanksmarkseries{arabic}
\renewenvironment{abstract}
  {\begin{center}\Large\sffamily\bfseries\abstractname\end{center}\begin{quotation}\sffamily\normalsize\setlength{\parindent}{0pt}}
  {\end{quotation}}
\fancypagestyle{plain}{%
  \fancyhf{}
  \fancyfoot[C]{\sffamily\bfseries\fontsize{11pt}{11pt}\selectfont\thepage}
  \renewcommand{\headrulewidth}{0pt}
  \renewcommand{\footrulewidth}{0pt}}
\pagestyle{plain}
\AtBeginEnvironment{thebibliography}{\small\raggedright\setlength{\parskip}{0pt}}
\numberwithin{equation}{section}
\renewcommand{\theequation}{\sffamily\thesection.\arabic{equation}}

\definecolor{gloryred}{HTML}{BB133E}
\definecolor{gloryblue}{HTML}{2C3ABC}
\definecolor{glorygreen}{HTML}{24802A}
\definecolor{topgrey}{HTML}{CFCFCF}
\definecolor{midgrey}{HTML}{AAAAAA}
\definecolor{downgrey}{HTML}{828282}
\newenvironment{statementbar}[1]
  {\def\FrameCommand{{\color{#1}\vrule width 3pt}\hspace{10pt}}%
   \MakeFramed{\advance\hsize-\width\FrameRestore}}
  {\endMakeFramed}
\declaretheoremstyle[
  headfont=\sffamily\bfseries,notefont=\sffamily\bfseries,
  notebraces={}{},headpunct=,bodyfont=\sffamily,
  headformat=\NAME~\NUMBER:~\NOTE\hfill\smallskip\linebreak,
  preheadhook=\begin{statementbar}{gloryred},postfoothook=\end{statementbar}
]{frozensTheorem}
\declaretheoremstyle[
  headfont=\sffamily\bfseries,notefont=\sffamily\bfseries,
  notebraces={}{},headpunct=,bodyfont=\sffamily,
  headformat=\NAME~\NUMBER:~\NOTE\hfill\smallskip\linebreak,
  preheadhook=\begin{statementbar}{gloryblue},postfoothook=\end{statementbar}
]{frozensLemma}
\declaretheoremstyle[
  headfont=\sffamily\bfseries,notefont=\sffamily\bfseries,
  notebraces={}{},headpunct=,bodyfont=\sffamily,
  headformat=\NAME~\NUMBER:~\NOTE\hfill\smallskip\linebreak,
  preheadhook=\begin{statementbar}{glorygreen},postfoothook=\end{statementbar}
]{frozensProposition}
\declaretheoremstyle[
  headfont=\sffamily\bfseries,notefont=\sffamily\bfseries,
  notebraces={}{},headpunct=,bodyfont=\sffamily,
  headformat=\NAME~\NUMBER:~\NOTE\hfill\smallskip\linebreak,
  preheadhook=\begin{statementbar}{midgrey},postfoothook=\end{statementbar}
]{frozensDefinition}
\declaretheoremstyle[
  headfont=\sffamily\bfseries,notefont=\sffamily\bfseries,
  notebraces={}{},headpunct=,bodyfont=\sffamily,
  headformat=\NAME~\NUMBER:~\NOTE\hfill\smallskip\linebreak,
  preheadhook=\begin{statementbar}{topgrey},postfoothook=\end{statementbar}
]{frozensRemark}
\declaretheorem[name=Theorem,style=frozensTheorem,numberwithin=section]{theorem}
\declaretheorem[name=Lemma,style=frozensLemma,sibling=theorem]{lemma}
\declaretheorem[name=Corollary,style=frozensTheorem,sibling=theorem]{corollary}
\declaretheorem[name=Proposition,style=frozensProposition,sibling=theorem]{proposition}
\declaretheorem[name=Hypothesis,style=frozensDefinition,sibling=theorem]{hypothesis}
\declaretheorem[name=Definition,style=frozensDefinition,sibling=theorem]{definition}
\declaretheorem[name=Remark,style=frozensRemark,sibling=theorem]{remark}
\DeclareMathOperator{\im}{im}
\DeclareMathOperator{\supp}{supp}
\DeclareMathOperator{\Mat}{Mat}
\DeclareMathOperator{\diag}{diag}
\DeclareMathOperator{\id}{id}
\newcommand{\NN}{\mathbb N}
\newcommand{\ZZ}{\mathbb Z}
\newcommand{\Prob}{\mathbb P}
\newcommand{\reliso}{\mathrel{\cong_{\Omega}}}
\newcommand{\Binf}{\mathsf B_\infty}
\newcommand{\cF}{\mathcal F}
\newcommand{\cY}{\mathcal Y}
\newcommand{\ev}{\operatorname{ev}}
\title{Finitary collapse of Bernoulli shifts\\from an algebraic splitting}
\author{Peter Burton\thanks{\href{mailto:peterburton1728@gmail.com}{\texttt{peterburton1728@gmail.com}}} - GPT-6 written draft.\thanks{Human-facing version to come soon; this manuscript has not been fully human-checked. Please contact author for more information.}}
\date{25 October 2026}
\begin{document}
\maketitle
\begin{abstract}
We construct a countable group for
which all nontrivial Bernoulli shifts
with countable discrete bases are finitarily conjugate, regardless of
their base entropies. For Polish bases the conjugacies are homeomorphisms
between invariant conull sets. The core construction turns a specified
finite-field rectangular splitting into reversible extraction of
countably many independent coordinates. The main external input is an algebraic existence result in Hypothesis~\ref{hyp:algebra} due to
OpenAI~\cite{OAI}. \\

Every original coordinate is
recovered after finitely many stages, and every auxiliary factor is
explicitly absorbed over one retained finite-base marker. We then combine
sparse initial data and a cheap marker with relative finitary integer
coding to treat all bases. The group can be taken to be the direct product
of the algebraic source group with the lamplighter group. Replacing the
lamplighter factor by an explicit finitely presented metabelian overgroup
gives a finitely presented example.
The core exhaustion argument is self-contained;
the all-base step uses two stated theorems of Meyerovitch and Spinka.\\

The external algebraic existence input is identified in
Appendix~\ref{sec:source} and is not reproved here.
An accompanying Lean certificate establishes the finite-uniform measurable
conjugacy described below; the all-base and finitary statements in this
note are beyond the scope of that certificate.
\end{abstract}

\section{Statement and proof plan}

\paragraph{Attribution and scope.}
This note develops a dynamical construction from the algebraic data of
Hypothesis~\ref{hyp:algebra}. The existence of those data is credited to
OpenAI's characteristic-two construction~\cite{OAI}; the pinned source and
the exact dependency are recorded in Appendix~\ref{sec:source}.
The purpose here is to explain how extraction, almost-sure recovery, and
absorption turn that input into Bernoulli conjugacies, and then to give the
stronger all-base finitary extension using the stated coding results of
Meyerovitch and Spinka~\cite{MS}.

\paragraph{Context as of 25 October 2026.}
For every countably infinite group, Bernoulli shifts with equal base
Shannon entropy are measurably isomorphic by Seward's theorem~\cite{Seward}.
The unequal-entropy question is different: Bowen's sofic entropy invariant
rules out a measurable conjugacy between finite-uniform Bernoulli shifts
of unequal base entropy over a sofic group~\cite{BowenSofic}.
Neither equal values of Rokhlin entropy nor factor maps in both directions
supply the inverse required for a conjugacy.

Related public approaches should also be acknowledged. The Astra research
archive develops finite defect extraction with a residual factor and a
conditional random equivariant basis criterion~\cite{Astra}. The
SauersML group-approximation project gives a conditional route through
approximate right inverses~\cite{Sauers}. In the versions consulted,
these approaches leave a residual factor, an unconstructed basis, or an
additional exactification hypothesis. They contain some antecedents to the
splitting and extraction strategy.

\paragraph{Formal verification and availability.}
The Lean build is available in the same GitHub repository as the compiled
PDF. Its closed finite-uniform certificate constructs a concrete countable
group and a finite alphabet with more than one symbol, and gives a
measure-preserving equivariant conjugacy on invariant conull sets between
the uniform Bernoulli shift on that alphabet and the shift on its square.
The algebraic witness is imported from the pinned OpenAI proofs. This
certifies the finite-uniform measurable existence statement; it does not
certify this note's all-base extension, the integer coding inputs,
finitary continuity, or the later corollaries. The stronger statements
are presented here with their written proofs and explicit external
dependencies.

All groups are countable and discrete. All products carry their product
Borel structures and the probability measures specified below. An
\emph{isomorphism of probability-preserving actions}, or \emph{conjugacy},
means an equivariant measure-preserving bijection between invariant conull
Borel sets, with measurable inverse.

If $(L,\lambda)$ is a probability space, its Bernoulli $G$-shift is
$(L^G,\lambda^G)$ with
\[
 (a\cdot x)(g)=x(a^{-1}g)\qquad(a,g\in G).
\]
The base Shannon entropy of the uniform measure on a finite set of size
$r$ is $\log r$. All logarithms are natural.

\begin{hypothesis}[The algebraic input]\label{hyp:algebra}
There are a finite field $K$, a countable group $N$, positive integers
$d,m,\kappa$, and subgroups $J_1,\ldots,J_d\leq N$ of order $\kappa$,
where $\kappa\cdot1_K\ne0$, with the following properties. Put
\[
 p_v=\kappa^{-1}\sum_{u\in J_v}u\in K[N],
 \qquad P=\diag(p_1,\ldots,p_d).
\]
There are matrices
\[
 U\in\Mat_{m\times d}(K[N]),\qquad
 T\in\Mat_{d\times m}(K[N])
\]
such that
\begin{equation}\label{eq:input}
 UP=U,\qquad PT=T,\qquad UT=I_m,\qquad d<\kappa m.
\end{equation}
Here $K[N]$ is the algebra of finite formal $K$-linear combinations of
elements of $N$, with multiplication induced by the group law.
\end{hypothesis}

Write $q=|K|$ and set
\begin{equation}\label{eq:parameters}
\begin{gathered}
 M=\kappa m,\qquad \Delta=M-d>0,
 \qquad A=\bigoplus_{j\in\ZZ}C_2,\\
 \Lambda=A\rtimes\langle s\rangle=C_2\wr\ZZ,\qquad
 D=N\times A,\qquad G=N\times\Lambda.
\end{gathered}
\end{equation}
where $s$ shifts the lamp indices by one. The height homomorphism
$G\to\ZZ$ kills $D$ and sends $s$ to $1$. In particular,
$D\cap\langle s\rangle=\{e\}$.
We regard $N$ and $A$ as the corresponding subgroups of $D$ and then of $G$,
and fix an enumeration of the lamp indices $\ZZ$.

\begin{definition}[Finitary regularity]\label{def:finitary}
For Polish probability bases, a conjugacy is \emph{topo-finitary} if it
is a homeomorphism between invariant conull Borel subsets of the product
spaces. For countable discrete bases this means that, at almost every
input, each output label is determined by finitely many input labels,
in both directions. We call this \emph{finitary}. The finite inspection
set can depend on the input; no uniform bound or moment bound is required.
For a nondiscrete output alphabet, continuity does not require an entire
output value to be determined by finitely many input labels.
\end{definition}

\begin{theorem}[All-base finitary collapse]\label{thm:all-bases}
Under Hypothesis~\ref{hyp:algebra}, any two nontrivial Bernoulli $G$-shifts
with Polish probability bases are topo-finitarily conjugate. In particular,
any two nontrivial countable-discrete Bernoulli bases give finitarily
conjugate $G$-shifts, including bases of unequal entropy.
\end{theorem}

Here nontrivial means that the probability measure is not concentrated at
one point. All nontrivial standard probability bases are covered
measurably, since their standard Borel spaces admit Polish realizations.
The first step is the following explicit finite-uniform construction.

\begin{theorem}\label{thm:main}
Under Hypothesis~\ref{hyp:algebra}, the following $G$-actions, with uniform
product measures, are measurably conjugate:
\begin{equation}\label{eq:main}
 (K^{d+1})^G\ \cong\ (K^{\NN_{>0}})^G.
\end{equation}
Consequently
\[
 (K^{d+1})^G\ \cong\ (K^{2(d+1)})^G.
\]
The latter two bases are finite and have distinct entropies
$(d+1)\log q$ and $2(d+1)\log q$.
\end{theorem}

The external construction in Appendix~\ref{sec:source} supplies
Hypothesis~\ref{hyp:algebra}. The group in this theorem is an enlargement
of its group $N$. We make no assertion here that the original group $N$
has the same conjugacy.

\paragraph{What must be proved.}
An algebraic splitting will let us repeatedly extract independent
coordinates. That alone does not ensure invertibility: information might
remain forever in the unextracted input. We therefore arrange that each
original coordinate stops depending on that input after finitely many
stages, almost surely. This is the main argument, in
Section~\ref{sec:exhaustion}.

Here is the structure of the proof. A single independent marker
$\Omega=K^G$ supplies fresh random orderings inside each left $D$-coset.
For every finite $H\leq A$, the algebra gives an exact splitting of $d$
coordinates per $H$-coset into $M>d$ such coordinates and an auxiliary
factor. At each stage we enlarge $H$, retain only one packet per larger
coset, and extract all other packets. We then apply enough algebraic
splittings to obtain $\Delta$ ordinary Bernoulli coordinates per group
element. The chance that any prescribed original coordinate still depends
on the retained input falls by a factor of at least two. After this
exhaustion, elementary countable-product isomorphisms absorb the auxiliary
factors.

\section{Coordinates, retained markers, and finite rows}\label{sec:notation}

For a finite subgroup $H\leq A$ and an integer $r\geq1$, define
\[
 X_H(r)=(K^r)^{G/H},\qquad
 X_H(\infty)=(K^{\NN_{>0}})^{G/H},\qquad
 \Binf=X_{\{e\}}(\infty).
\]
The set $G/H$ consists of the left cosets $gH$. A scalar coordinate of
$X_H(r)$ is indexed by $(gH,i)$, $1\leq i\leq r$. These are compact
metrizable abelian groups under addition, with their Haar probability
measures, namely independent uniform $K$-labels. The left $G$-action
permutes the cosets.

We use the word \emph{packet} for the $r$ scalar labels at one coset.
Every coset $gH$ is contained in a unique left $D$-coset. A map
\emph{preserves $D$-cosets} if its output on one such coset uses only
data in that coset. Marker digits themselves may lie elsewhere; the
random variables indexed by this $D$-coset will be specified explicitly.

\begin{definition}[Isomorphisms over a marker]\label{def:relative}
Let $\Omega=K^G$, with uniform product measure $\Prob$, and let $X,Y$ be
probability-preserving $G$-spaces independent of $\Omega$. The notation
$X\reliso Y$ means that there is a conjugacy
\[
 (\omega,x)\longmapsto(\omega,\phi_\omega(x))
 \quad\text{from }\Omega\times X\text{ to }\Omega\times Y,
\]
where the map and inverse are jointly measurable, and, for almost every
$\omega$, $\phi_\omega$ carries the probability measure on $X$ to that on
$Y$. Equivariance means
\begin{equation}\label{eq:covariance}
 \phi_{a\omega}(ax)=a\phi_\omega(x)\qquad(a\in G).
\end{equation}
The same single marker occurs on both sides.
\end{definition}

\begin{lemma}[Products over the same marker]\label{lem:relative-products}
Suppose $X_j\reliso Y_j$ for $j\geq1$. Applying all the fiber maps
coordinatewise gives $\prod_j X_j\reliso\prod_j Y_j$. These relative
isomorphisms can also be composed and inverted.
\end{lemma}
\begin{proof}
Intersect the countably many invariant conull sets of allowed markers.
For each marker in this intersection, the product map is a bijection
with its coordinatewise inverse. Both are measurable because all
coordinate maps are measurable. On cylinder sets the product measure is
preserved, hence it is preserved on the product Borel sigma-algebra.
Equation~\eqref{eq:covariance} holds coordinatewise. Integrating in
$\omega$ proves preservation of the full product measure. Composition
and inversion have the same properties.
\end{proof}

\begin{definition}[Rows and inverse sections]\label{def:rows}
A \emph{row} on $X_H(r)$ is a finitely supported linear functional
\[
 \ell(x)=\sum_{(b,i)\in F}c_{b,i}x_i(b),
 \qquad F\subseteq (G/H)\times\{1,\ldots,r\}\text{ finite}.
\]
Its support consists of the nonzero coefficients, after collecting terms
with the same index. A linear map has row bound $C$ if every output
scalar coordinate is a row with at most $C$ terms. Composing maps of
row bounds $C_1,C_2$ gives row bound at most $C_1C_2$.

If $\phi:X\to X'\times Y$ is a linear isomorphism, its \emph{inverse
section for $X'$} is $x'\mapsto\phi^{-1}(x',0)$. It describes the part
of the old input depending on the new input when the extracted factor
$Y$ is set to zero. It is not asserted to preserve probability measure.
\end{definition}

\section{Fresh randomness from one finite-base marker}\label{sec:marker}

Fix a bijection $\beta:K\to\{0,\ldots,q-1\}$ and a bijection
\[
 \rho:\NN_{\geq0}\times\NN_{>0}\longrightarrow\NN_{>0}.
\]
For $g\in G$ and a stream index $r\geq0$, put
\begin{equation}\label{eq:uniforms}
 \xi_g^{(r)}(\omega)
   =\sum_{j=1}^{\infty}q^{-j}
          \beta\bigl(\omega(gs^{\rho(r,j)})\bigr).
\end{equation}
These are measurable $[0,1]$-valued functions of the marker.

\begin{lemma}\label{lem:marker}
For each fixed left $D$-coset $C$, the family
$\{\xi_g^{(r)}:g\in C,\ r\geq0\}$ consists of independent uniform
random variables on $[0,1]$. Moreover
\[
 \xi_{ag}^{(r)}(a\omega)=\xi_g^{(r)}(\omega).
\]
\end{lemma}
\begin{proof}
Write $C=g_0D$. The elements $g_0us^j$, for $u\in D$ and $j\geq1$,
are pairwise distinct. Indeed, equality for $(u,j)$ and $(v,k)$ would
give $v^{-1}u=s^{k-j}$. Applying the height homomorphism gives
$j=k$, and then $u=v$.
The digit sequences used in \eqref{eq:uniforms}, for all the indicated
$(g,r)$, are disjoint sequences of independent uniform digits. A sequence
of such base-$q$ digits has the uniform distribution on $[0,1]$: each
interval specified by its first $j$ digits has probability $q^{-j}$,
and these intervals determine the distribution. This also proves
independence of the family. The covariance identity follows directly
from $(a\omega)(ag s^j)=\omega(gs^j)$.
\end{proof}

For a finite $H\leq A$, give an $H$-coset $b$ the priority
\begin{equation}\label{eq:priorities}
 \pi_{H,r}(b)=\min_{g\in b}\xi_g^{(r)}.
\end{equation}
Distinct $H$-cosets inside a fixed $D$-coset have independent priorities
with the same continuous distribution. A finite collection of these
cosets can therefore be ordered by increasing priority, and every member
is equally likely to come first. Ordering the individual elements of a
finite coset uses $H=\{e\}$.

There are only countably many finite subgroups of $A$, cosets, stream
indices, and pairs of priorities being compared. The event on which all
these comparisons are tie-free has probability one and is $G$-invariant.
Call it $\Omega_{\mathrm{ord}}$. For measurability statements on the
whole probability space, break any ties using a fixed enumeration of the
countable index sets. This convention has no effect on
$\Omega_{\mathrm{ord}}$, where every ordering is covariant.

\begin{remark}\label{rem:local-independence}
The random fields belonging to different left $D$-cosets need not be
independent. We will never require this. The recovery estimate for a
coordinate conditions only on earlier streams indexed by its own
$D$-coset. We also do not condition on the entire marker in that estimate:
once the marker is fixed, all the selections are fixed.
\end{remark}

\section{The algebraic splitting on each finite coset space}\label{sec:split}

\subsection{The action of the group algebra}

For $a=\sum_{u\in N}a_u u\in K[N]$, define an operator on $K^{G/H}$ by
\begin{equation}\label{eq:operator}
 (\mathcal R_a f)(gH)=\sum_{u\in N}a_u f(guH).
\end{equation}
This is well-defined: replacing $g$ by $gh$, $h\in H$, does not change
$guH$, since $H\leq A$ commutes with $N$. Each sum is finite. Direct
substitution gives
\[
 \mathcal R_a\mathcal R_b=\mathcal R_{ab}.
\]
For a matrix $B$ over $K[N]$, use
$(\mathcal R_B f)_i=\sum_j\mathcal R_{B_{ij}}f_j$.
These continuous linear operators commute with the left $G$-action,
preserve $D$-cosets, and respect matrix multiplication. We henceforth
write the matrices themselves for their operators on $G/H$.

For every $v$, $p_v^2=p_v$. Its image consists exactly of the functions
constant on each coset of $HJ_v$. The inclusion in this space follows
by reindexing the subgroup average, and a function already constant
there is fixed by that average. Because $H\cap J_v=\{e\}$ and the two
subgroups commute, each $HJ_v$-coset contains exactly $\kappa$
$H$-cosets. It follows that
\begin{equation}\label{eq:VH}
 V_H:=\im(P:X_H(d)\to X_H(d))
      =\prod_{v=1}^d K^{G/(HJ_v)}.
\end{equation}
This identification preserves Haar probability measure.

\begin{lemma}[A split product]\label{lem:algebraic-split}
Put $Q=P-TU$ and let $Y_H=\im(Q:X_H(d)\to X_H(d))$. Then
\[
 V_H\longrightarrow X_H(m)\times Y_H,
 \qquad v\longmapsto(Uv,Qv)
\]
is a continuous linear $G$-equivariant isomorphism, with inverse
$(y,z)\mapsto Ty+z$. It sends Haar measure to product Haar measure.
\end{lemma}
\begin{proof}
The identities in \eqref{eq:input} give
\[
 Q^2=Q,\quad PQ=QP=Q,\quad UQ=0,\quad QT=0.
\]
For example, $(TU)^2=TU$, while $PTU=TUP=TU$, so
$(P-TU)^2=P-TU$ in every characteristic.
If $v\in V_H$, then $TUv+Qv=Pv=v$.
Conversely, for $y\in X_H(m)$ and $z\in Y_H$, $Ty+z\in V_H$ and
\[
 U(Ty+z)=y,\qquad Q(Ty+z)=z.
\]
All maps are continuous because their coordinate formulas are finite.
The image $Y_H$ is compact, hence closed. A continuous group isomorphism
between compact groups carries normalized Haar measure to normalized
Haar measure; Haar measure on a product is the product of the Haar
measures. This proves the measure assertion.
\end{proof}

\subsection{Turning subgroup-constant labels into ordinary coset labels}

Take $\kappa$ independent copies of the splitting in
Lemma~\ref{lem:algebraic-split}. Define the fixed compact $G$-space
\[
 Z_H=Y_H^{\kappa},
\]
with product Haar measure. Fix a stream index $r$. In each $HJ_v$-coset
$b$, order its $\kappa$ constituent $H$-cosets as
$b_1,\ldots,b_\kappa$ by \eqref{eq:priorities}. Define
\begin{equation}\label{eq:F}
 F_{H,r,\omega}:V_H^\kappa\longrightarrow X_H(d),
 \qquad [F_{H,r,\omega}(v)]_v(b_j)=v_v^{(j)}(b).
\end{equation}
Thus the label in copy $j$ is placed at the $j$th ordered coset.
For a fixed marker this is a bijection of scalar coordinate slots,
hence a continuous linear Haar-preserving isomorphism. Its map and
inverse are jointly measurable in the marker, and they satisfy
\eqref{eq:covariance} on $\Omega_{\mathrm{ord}}$.

Combining $F_{H,r,\omega}^{-1}$ with the $\kappa$ copies of the algebraic
split, and identifying $(X_H(m))^\kappa$ with $X_H(M)$, gives
\begin{equation}\label{eq:relative-split}
 X_H(d)\reliso X_H(M)\times Z_H.
\end{equation}
Fix the track order in this identification once and for all. Separating
the first $d$ tracks gives
\begin{equation}\label{eq:one-step}
 X_H(d)\reliso X_H(d)\times X_H(\Delta)\times Z_H.
\end{equation}
Conditional on the marker, all factors on the right have their product
Haar law. In particular, the new $X_H(d)$ input is independent of the
extracted factors. This statement concerns the full split map; its
inverse section alone does not have this measure property.

\subsection{A uniform finite bound on inverse rows}

Define the finite integer
\begin{equation}\label{eq:C}
 C_0=\max\left(1,\max_{1\leq v\leq d}
                    \sum_{i=1}^m|\supp(T_{vi})|\right).
\end{equation}
Let $\widehat T$ apply $T$ separately in the $\kappa$ copies. The inverse
section for the $X_H(M)$ factor in \eqref{eq:relative-split} is
\[
 B_{H,r,\omega}=F_{H,r,\omega}\widehat T.
\]
Each row of this map uses at most $C_0$ scalar coordinates: $F$ chooses
one coordinate in one copy, and that row of $T$ has at most $C_0$ terms.
Collisions of cosets or cancellation can only reduce this bound. It is
uniform in $H,r,\omega$.

Let $\iota_d:X_H(d)\to X_H(M)$ insert the input into the first $d$
tracks, with zero in the other tracks. Then
\begin{equation}\label{eq:S}
 S_{H,r,\omega}=B_{H,r,\omega}\iota_d
\end{equation}
is the inverse section for the first factor of \eqref{eq:one-step}.
Iterate that same split $k$ times, each time on its first factor. We
obtain
\begin{equation}\label{eq:k-split}
 X_H(d)\reliso
 X_H(d)\times X_H(k\Delta)\times Z_H^k.
\end{equation}
For a fixed marker its inverse section is $S_{H,r,\omega}^{\,k}$,
with row bound $C_0^k$. All these maps preserve $D$-cosets, and their
coefficients on a given $D$-coset depend only on stream $r$ indexed
by that coset. The same stream may be reused in all $k$ iterations:
conditional product Haar measure is preserved at each iteration.

\section{Growing finite blocks and recovery of every coordinate}\label{sec:exhaustion}

\subsection{Choose all sizes deterministically}

Let $A_a$ denote the subgroup of $A$ generated by the first $a$ lamps
in the fixed enumeration, so $|A_a|=2^a$ and $A_0=\{e\}$. Define increasing integers
$a_n$, finite subgroups $H_n=A_{a_n}$, and row bounds $R_n$ recursively.
Start with $a_0=0$ and $R_0=1$. Having chosen $a_n$ and $R_n$, choose
$a_{n+1}>a_n$ so that
\begin{equation}\label{eq:choose-sizes}
 L_n:=2^{a_{n+1}-a_n}=[H_{n+1}:H_n]\geq2R_n.
\end{equation}
Then set
\begin{equation}\label{eq:recursion}
 h_{n+1}=|H_{n+1}|,
 \qquad k_n=h_{n+1},
 \qquad R_{n+1}=R_n C_0^{k_n}.
\end{equation}
Every choice is finite. The order is important: $R_n$ determines a
sufficient block size, that size determines $k_n$, and only then do we
define $R_{n+1}$. There is no requirement to bound $R_{n+1}$ before
choosing the current block.

\subsection{One stage of the construction}

At stage $n$, the remaining input is $x_n\in X_{H_n}(d)$. Each
$H_{n+1}$-coset $B$ is a union of $L_n$ cosets of $H_n$. Order them
using stream $2n$:
\[
 b_1(B),\ldots,b_{L_n}(B).
\]
Retain the $d$-packet at $b_1(B)$ and extract the other packets in their
listed order. This is an isomorphism over $\Omega$
\begin{equation}\label{eq:regroup}
 X_{H_n}(d)\reliso X_{H_{n+1}}(d)\times W_n,
 \qquad W_n=X_{H_{n+1}}\bigl(d(L_n-1)\bigr).
\end{equation}
The inverse section for its first factor is the insertion map
\begin{equation}\label{eq:insertion}
 (I_{n,\omega}v)_i(b)=
 \begin{cases}
 v_i(B),& b=b_1(B),\quad b\subset B,\\
 0,& b\ne b_1(B),\quad b\subset B.
 \end{cases}
\end{equation}
Each row of $I_{n,\omega}$ has at most one nonzero coefficient.

Next apply \eqref{eq:k-split} to the retained factor, with
$H=H_{n+1}$, stream $2n+1$, and $k=k_n$. The complete stage is
\begin{equation}\label{eq:stage}
 X_{H_n}(d)\reliso X_{H_{n+1}}(d)\times\cY_n,
 \qquad
 \cY_n=W_n\times E_n\times Z_{H_{n+1}}^{k_n},
\end{equation}
where
\begin{equation}\label{eq:En}
 E_n=X_{H_{n+1}}(k_n\Delta)
     =X_{H_{n+1}}(h_{n+1}\Delta).
\end{equation}
Write its fiber map as
$\phi_{n,\omega}(x_n)=(x_{n+1},y_n)$, with $y_n\in\cY_n$.
Its inverse section for the new input is
\begin{equation}\label{eq:stage-section}
 J_{n,\omega}
   =I_{n,\omega}S_{H_{n+1},2n+1,\omega}^{\,k_n}.
\end{equation}
Composition is read from right to left: the algebraic section acts
first; insertion into the one selected old packet acts afterwards.
This map has row bound $C_0^{k_n}$.

\begin{lemma}[A fixed Bernoulli output at every stage]\label{lem:ordinary-output}
For every $n$, $E_n\reliso X_{\{e\}}(\Delta)$.
\end{lemma}
\begin{proof}
In each $H_{n+1}$-coset $B$, order its $h_{n+1}$ individual group
elements by their stream-$2n$ values. Regard the $h_{n+1}\Delta$ labels
at $B$ as $h_{n+1}$ ordered packets of $\Delta$ labels, and place packet
$j$ at the $j$th element. This is a measurable covariant bijection of
scalar slots. For every tie-free marker it is a linear homeomorphism
and preserves uniform product measure.
\end{proof}

The reuse of stream $2n$ here causes no difficulty: for a fixed marker
this operation is just a coordinate permutation. Since every full
stage preserves conditional product Haar measure, the extracted
ordinary Bernoulli factors are independent of one another, the other
extracted factors, and the retained marker.

\subsection{The dependence of an original coordinate}

Let $\Pi_{0,\omega}$ be the identity on $X_{H_0}(d)$, and recursively put
\[
 \Pi_{n+1,\omega}=\Pi_{n,\omega}J_{n,\omega}.
\]
Thus $\Pi_{n,\omega}:X_{H_n}(d)\to X_{H_0}(d)$ is the inverse section
for the input left after the first $n$ stages, with all extracted
factors set to zero. For an original scalar coordinate
$\alpha=(g,i)\in G\times\{1,\ldots,d\}$, define the row
\begin{equation}\label{eq:ell}
 \ell_{n,\alpha,\omega}=\ev_\alpha\Pi_{n,\omega}.
\end{equation}
The row is zero exactly when that original coordinate no longer depends
on the remaining input.

Fix this original coordinate and let $C=gD$. Define
\begin{equation}\label{eq:filtration}
 \cF_n^C=\sigma\bigl(\xi_u^{(r)}:
                          u\in C,\ 0\leq r<2n\bigr).
\end{equation}
All probability statements in the next lemma are on the original
marker space $(\Omega,\Prob)$, before restricting to a conull set.

\begin{lemma}[Recovery estimate]\label{lem:recovery}
The row $\ell_{n,\alpha,\omega}$ has at most $R_n$ nonzero scalar
coefficients. It is supported on $H_n$-cosets inside $C$, and its
coefficients are $\cF_n^C$-measurable. Moreover
\begin{equation}\label{eq:conditional-survival}
 \Prob\bigl(\ell_{n+1,\alpha,\omega}\ne0\mid\cF_n^C\bigr)
 \leq \frac{R_n}{L_n}\,
              \mathbf1_{\{\ell_{n,\alpha,\omega}\ne0\}}
 \leq \frac12\,
              \mathbf1_{\{\ell_{n,\alpha,\omega}\ne0\}}
 \quad\text{almost surely}.
\end{equation}
\end{lemma}
\begin{proof}
The assertions on support and measurability follow by induction.
At $n=0$ we have the fixed coordinate evaluation, of support size one.
The section $J_{n,\omega}$ has row bound $C_0^{k_n}$, preserves
$D$-cosets, and uses only streams $2n$ and $2n+1$ within the relevant
coset. The recursion \eqref{eq:recursion} therefore gives all three
assertions at $n+1$. Measurability of a row means measurability of its
coefficient at each of the countably many possible coordinate indices.

Condition now on $\cF_n^C$. The row $\ell_n=\ell_{n,\alpha,\omega}$
is fixed. Let $\mathcal S_n$ be the set of $H_n$-cosets on which at
least one of its coefficients is nonzero. Then
$|\mathcal S_n|\leq R_n$. By Lemma~\ref{lem:marker}, the stream-$2n$
priorities within $C$ are independent of $\cF_n^C$. In every
$H_{n+1}$-coset, the $L_n$ constituent $H_n$-cosets have independent,
identically distributed continuous priorities. Thus each specified
$H_n$-coset is selected with conditional probability $1/L_n$.

If no member of $\mathcal S_n$ is selected, then
$\ell_n I_{n,\omega}=0$ by \eqref{eq:insertion}. In that event
\[
 \ell_{n+1,\alpha,\omega}
   =\ell_n I_{n,\omega}
          S_{H_{n+1},2n+1,\omega}^{\,k_n}=0,
\]
regardless of the second stream or the algebraic spreading. A
conditional union bound gives $|\mathcal S_n|/L_n$ as an upper bound
on survival. If $\ell_n=0$, its successor is also zero. Combining these
facts proves \eqref{eq:conditional-survival}. No independence between
the selections in different blocks is needed.
\end{proof}

Taking expectations and iterating gives
\begin{equation}\label{eq:survival}
 \Prob(\ell_{n,\alpha,\omega}\ne0)\leq2^{-n}.
\end{equation}
The events on the left are decreasing in $n$, because a zero row stays
zero. Continuity of probability for decreasing events therefore yields
\[
 \Prob(\ell_{n,\alpha,\omega}\ne0\text{ for every }n)=0.
\]
There are countably many original scalar coordinates. Consequently
there is a Borel conull set $\Omega_*\subseteq\Omega_{\mathrm{ord}}$
on which \emph{every} original coordinate row becomes zero at a finite
stage. All sections are covariant, so the property quantified over all
coordinates is invariant under $G$. Thus $\Omega_*$ is $G$-invariant.

\subsection{The infinite extraction has a measurable inverse}

\begin{proposition}\label{prop:exhaustion-isomorphism}
The successive stage maps give an isomorphism over the marker
\begin{equation}\label{eq:infinite-extraction}
 X_{H_0}(d)\reliso\prod_{n\geq0}\cY_n.
\end{equation}
For every $\omega\in\Omega_*$, the fiber map is a continuous linear
isomorphism of compact groups.
\end{proposition}
\begin{proof}
Fix $\omega\in\Omega_*$. Starting from $x_0$, recursively apply
$\phi_{n,\omega}$ to obtain $x_{n+1}$ and $y_n$. Define
\[
 \mathcal E_\omega(x_0)=(y_0,y_1,\ldots).
\]
This map is linear and continuous: each output coordinate is obtained
by finitely many continuous stage maps. For every $n$, composing the
first $n$ stages gives a continuous linear isomorphism
\begin{equation}\label{eq:finite-prefix}
 X_{H_0}(d)\longrightarrow
 X_{H_n}(d)\times\prod_{j<n}\cY_j.
\end{equation}

\emph{Injectivity.} The inverse of \eqref{eq:finite-prefix} expresses
each original coordinate as
\begin{equation}\label{eq:coordinate-recovery}
 x_0(\alpha)=
 f_{n,\alpha,\omega}(y_0,\ldots,y_{n-1})
       +\ell_{n,\alpha,\omega}(x_n),
\end{equation}
where $f_{n,\alpha,\omega}$ is a continuous linear function of the
extracted data. For each $\alpha$ choose a stage with
$\ell_{n,\alpha,\omega}=0$. That coordinate is then determined by
the extracted data alone. Hence two inputs with the same infinite
extraction agree at every coordinate.

\emph{Surjectivity.} The image of $\mathcal E_\omega$ is compact and
therefore closed in $\prod_n\cY_n$. Its projection onto every finite
initial product $\prod_{j<n}\cY_j$ is onto, by
\eqref{eq:finite-prefix}. Every basic open set in the infinite product
restricts only finitely many factors and is therefore met by the image.
The image is dense as well as closed, so it is the whole product.
This proves that $\mathcal E_\omega$ is a compact-group isomorphism,
and it consequently preserves Haar probability measure.

\emph{Joint measurability of the inverse.} Each coefficient of
$\ell_{n,\alpha,\omega}$ is measurable in $\omega$. Equality of this
row to zero is a countable intersection of coefficient equalities,
and is therefore a Borel event. On $\Omega_*$ define
\[
 \tau_\alpha(\omega)
   =\min\{n\geq0:\ell_{n,\alpha,\omega}=0\}.
\]
This is a finite Borel integer-valued function. The inverse coordinate
is given by the first term of \eqref{eq:coordinate-recovery}, with
$n=\tau_\alpha(\omega)$. Each fixed-$n$ expression is jointly
measurable, being part of a finite composition of the stage inverses.
Partitioning according to the countably many values of $\tau_\alpha$
proves joint measurability of this coordinate, and hence of the full
inverse. The forward map is jointly measurable for the same
finite-stage reason.

Finally all stage maps are covariant, so their infinite extraction and
its inverse are covariant on $\Omega_*$. Fiberwise preservation of
Haar measure, integrated against $\Prob$, proves the required
measure-preserving conjugacy.
\end{proof}

\section{Absorbing the remaining factors}\label{sec:absorption}

By Lemma~\ref{lem:ordinary-output}, the product of the factors $E_n$
in \eqref{eq:infinite-extraction} is, over the same marker, the ordinary
Bernoulli shift $\Binf$: there are countably many sets of $\Delta>0$
tracks, and a fixed bijection of their indices with $\NN_{>0}$
interleaves them. Thus
\begin{equation}\label{eq:with-residuals}
 X_{H_0}(d)\reliso
 \Binf\times\prod_{n\geq0}
       \bigl(W_n\times Z_{H_{n+1}}^{k_n}\bigr).
\end{equation}
We now remove the factors on the right by explicit isomorphisms.

\begin{lemma}[Finite-coset coordinate interleaving]\label{lem:interleave}
For every finite $H\leq A$,
\[
 X_H(\infty)\reliso\Binf.
\]
Consequently $\Binf\times X_H(r)\reliso\Binf$ for every finite $r$.
\end{lemma}
\begin{proof}
Write $h=|H|$. In each $H$-coset $b$, order its elements as
$g_1,\ldots,g_h$ using stream $0$. For $x\in X_H(\infty)$, define
\[
 y_k(g_i)=x_{h(k-1)+i}(b)
 \qquad(k\geq1,\ 1\leq i\leq h).
\]
This is a covariant permutation of scalar coordinate slots, with its
explicit inverse given by the same formula read backwards. It is a
linear Haar-preserving homeomorphism for each allowed marker and is
jointly measurable in both directions. On $X_H(\infty)$, one can
interleave any additional $r$ tracks with the countably many existing
ones using a fixed bijection $\NN_{>0}\sqcup\{1,\ldots,r\}\to\NN_{>0}$.
Conjugating that operation by the first isomorphism proves the second.
\end{proof}

\begin{lemma}[Absorption of an algebraic remainder]\label{lem:absorb-Z}
For every finite $H\leq A$ and every finite integer $k\geq1$,
\[
 \Binf\times Z_H^k\reliso\Binf.
\]
\end{lemma}
\begin{proof}
Take countably many copies of \eqref{eq:relative-split}, using the same
marker and, for example, stream $0$ for every copy. By
Lemma~\ref{lem:relative-products} this gives
\[
 (X_H(d))^{\NN_{>0}}\reliso
 (X_H(M))^{\NN_{>0}}\times Z_H^{\NN_{>0}}.
\]
Since $d,M\geq1$, fixed track interleavings identify each of the first
two countable products with $X_H(\infty)$. Lemma~\ref{lem:interleave}
therefore supplies an isomorphism
\begin{equation}\label{eq:absorb-countable}
 \theta_H:\Binf\reliso\Binf\times Z_H^{\NN_{>0}}.
\end{equation}
To absorb $k$ more copies, use the following composition, throughout
over the same marker:
\[
 \Binf\times Z_H^k
 \xrightarrow{\theta_H\times\id}
 \Binf\times Z_H^{\NN_{>0}}\times Z_H^k
 \xrightarrow{\text{interleave}}
 \Binf\times Z_H^{\NN_{>0}}
 \xrightarrow{\theta_H^{-1}}\Binf.
\]
The middle arrow is a permutation of the copies of the \emph{same}
probability-preserving action $Z_H$. This is an actual isomorphism;
no cancellation of factors is being inferred.
\end{proof}

\begin{proof}[Completion of Theorem~\ref{thm:main}]
First split the countably many scalar tracks of $\Binf$ into countably
many infinite sets. This identifies $\Binf$ with a countable product
of independent copies of itself. Allocate one copy to each $W_n$ and
one to each $Z_{H_{n+1}}^{k_n}$. Lemma~\ref{lem:interleave} absorbs
$W_n=X_{H_{n+1}}(d(L_n-1))$ into its allocated copy, and
Lemma~\ref{lem:absorb-Z} absorbs $Z_{H_{n+1}}^{k_n}$ into the other
allocated copy. Taking the product of these isomorphisms over the
single marker, then interleaving the resulting countable family of
$\Binf$ factors, yields
\begin{equation}\label{eq:all-absorption}
 \Binf\times\prod_{n\geq0}
       \bigl(W_n\times Z_{H_{n+1}}^{k_n}\bigr)
       \reliso\Binf.
\end{equation}
Combining \eqref{eq:with-residuals} and \eqref{eq:all-absorption} proves
$X_{H_0}(d)\reliso\Binf$. Restoring the retained marker gives
\[
 (K^{d+1})^G
     =\Omega\times X_{H_0}(d)
     \cong\Omega\times\Binf
     \cong\Binf.
\]
The final arrow interleaves the single marker coordinate with the
countably many data coordinates at each group element. This proves
\eqref{eq:main}.

Let $\Phi:(K^{d+1})^G\to\Binf$ denote this conjugacy. Then
\[
 (K^{2(d+1)})^G
 \cong (K^{d+1})^G\times(K^{d+1})^G
 \xrightarrow{\Phi\times\Phi}\Binf\times\Binf
 \xrightarrow{\text{interleave}}\Binf
 \xrightarrow{\Phi^{-1}}(K^{d+1})^G
\]
is the required conjugacy of two finite-base shifts. Since $q\geq2$
and $d\geq1$, their stated Shannon entropies are positive and unequal.
\end{proof}

\section{Sparse initial data and finitary regularity}\label{sec:sparse-finitary}

The construction above has two useful freedoms: its initial finite
subgroup need not be trivial, and its marker need not be uniform.
We record both, including the joint continuity needed below.

Fix $0<p<1$, let $\nu_p$ assign probabilities $p,1-p$ to $0,1$, and
write $\Omega_p=(\{0,1\}^G,\nu_p^G)$. In this section relative
isomorphisms retain $\Omega_p$ instead of the uniform marker.
Set
\[
 \epsilon_p=-p\log p-(1-p)\log(1-p).
\]

\begin{lemma}[Sparse-start relative exhaustion]\label{lem:sparse-start}
For every $a\geq0$ and $H=A_a$, there is a topo-finitary conjugacy
\[
 \Omega_p\times X_H(d)\ \cong\ \Omega_p\times\Binf
\]
which fixes the marker. For each allowed marker its data map is a
linear homeomorphism of compact groups preserving Haar measure.
\end{lemma}
\begin{proof}
First supply the random streams. In each disjoint digit column used in
Section~\ref{sec:marker}, group successive biased bits into pairs.
Discard $00$ and $11$, and read $01$ and $10$ as fair bits $0$ and $1$.
An accepted pair occurs with probability $2p(1-p)>0$, independently
between pairs. Almost surely infinitely many are accepted. Conditional
on their positions, the accepted bits are independent and fair, so
they remain so after forgetting the positions. Their binary expansion
is uniform on $[0,1]$. Disjoint input columns give independent streams
inside each fixed $D$-coset. Covariance is unchanged.

Every prescribed finite prefix of accepted bits requires finitely
many original bits. The binary-expansion map is continuous.
Thus a finite collection of strict priority comparisons is certified
by finitely many marker coordinates. Intersect the countably many
events of infinite acceptance and absence of ties to obtain an
invariant conull marker set.

Start the recursion of Section~\ref{sec:exhaustion} at $a_0=a$,
$H_0=H$, $R_0=1$. An original coordinate is now $(gH,i)$.
Its initial row still has one coefficient, supported in $gD$.
Choose $L_n\geq2R_n$ exactly as before and keep
$k_n=|H_{n+1}|$. This is the absolute subgroup order, not the
relative index $[H_{n+1}:H_n]$. Therefore
Lemma~\ref{lem:ordinary-output} still extracts $\Delta$ ordinary
$G$-tracks at each stage. All row estimates, compactness,
measurability, and absorption arguments apply without another change.
They give the asserted relative measurable compact-group isomorphism.

It remains to check joint continuity, including in the marker.
For a fixed finite stage and output scalar, expand its formula into
its full finite computation tree. Each algebraic operator has finite
support; each ordering concerns a finite coset. At the given marker,
finitely many strict comparisons determine every coordinate index in
this tree. Their finite certificates freeze all its branches. The
entire linear row, before collecting terms or cancellations, is then
fixed. This proves continuity of every finite-stage coordinate map,
both forward and backward, on the tie-free domain.

For an inverse coordinate $\alpha$, choose any stage $n$ at which
$\ell_{n,\alpha,\omega}=0$. Freeze the finite tree of this whole row,
not just the terms which survive cancellation. The same row is zero
for every allowed marker in that neighborhood. Freeze also the finite
tree of the recovery function $f_{n,\alpha,\omega}$ in
\eqref{eq:coordinate-recovery}. On this neighborhood the inverse
coordinate is one fixed finite linear function of extracted data.
It is therefore continuous. The first vanishing stage need not be
locally constant: the single chosen witness stage suffices.

Finally, the absorption maps of Section~\ref{sec:absorption} are
coordinate permutations and countable products of finite compositions
of elementary splits and their inverses. Each requested coordinate
uses only one such finite composition. There is no infinite recursive
application along one coordinate. Consequently these maps and their
inverses are jointly continuous on the common invariant conull
domain. Their composition with extraction proves the lemma.
\end{proof}

\section{From sparse data to every Bernoulli base}\label{sec:all-bases}

We now use two precise inputs from integer finitary coding.
For a probability base $\lambda$, write $H(\lambda)$ for its Shannon
entropy, equivalently the supremum of the entropies of its finite
measurable partitions. This is the sum of $-t\log t$ over its atoms
when it is purely atomic, and is infinite if it has a nonatomic part.

\begin{lemma}[Integer coding inputs]\label{lem:integer-coding}
The following consequences of Meyerovitch--Spinka~\cite{MS} will be used.
\begin{enumerate}[label=(\roman*)]
\item Equal-entropy iid integer processes with Polish bases admit a
topo-finitary conjugacy, also when their entropy is infinite.
\item Let $W$ be iid Bernoulli-$\theta$ on $\mathbb Z$, $0<\theta<1$.
Given $W$, let $X_n$ independently be uniform on $K^d$ when $W_n=1$
and the fixed symbol $*$ otherwise. If a finite probability base
$\lambda$ has entropy $\theta d\log q$, then $(W,X)$ is finitarily
isomorphic over $W$ to $W$ times an independent iid $\lambda$ process.
\end{enumerate}
\end{lemma}
\begin{proof}[Identification of the cited inputs]
Part (i) is \cite[Theorem 3.15]{MS}.
For (ii), apply \cite[Theorem 3.2]{MS}. Both compared processes are
conditionally independent over the aperiodic ergodic process $W$,
with relative entropy $\theta d\log q$. Their conditional laws on
any finite set are determined by the corresponding finite $W$-block.
Thus the theorem's relative dependence bound is zero and its
stop-finitary conditional-law hypothesis holds. Its relative
isomorphism fixes $W$. We use the resulting bijection, not only a
factor map.
\end{proof}

\begin{remark}[Keeping the topology of a Polish base]
The topological reduction underlying part (i) can be made explicit.
For a Polish probability space $(L,\lambda)$, choose a countable
metric-ball basis $(U_j)$ with $\lambda(\partial U_j)=0$.
This is possible because, at each center in a countable dense set,
only countably many spheres have positive measure.
The set $L_0=L\setminus\bigcup_j\partial U_j$ is conull and $G_\delta$.
Each $U_j\cap L_0$ is clopen in $L_0$. Consequently
$x\mapsto(\mathbf1_{U_j}(x))_j$ is a homeomorphism of $L_0$ onto
its image in $\{0,1\}^{\NN}$: the basis separates points and also
makes the inverse continuous. This realizes a conull part of the
original topology without replacing it by a discrete topology.
The integer coding assertion itself remains the cited theorem.
\end{remark}

The next elementary observation lets us apply integer coding along
the left cosets $g\langle s\rangle$, parametrized by right multiplication
by powers of $s$.

\begin{lemma}[Extension from a subgroup]\label{lem:coinduce}
If two Bernoulli shifts over a countable subgroup $B\leq\Gamma$
are topo-finitarily conjugate, their Bernoulli shifts over $\Gamma$
are topo-finitarily conjugate. Finitary conjugacy of discrete bases
is preserved as well.
\end{lemma}
\begin{proof}
Apply the same $B$-equivariant code on every left coset $gB$,
identified with $B$ by $b\mapsto gb$. Changing this identification
by a left translation in $B$ does not change the output, by
$B$-equivariance. Left translation by $\Gamma$ carries these
cosetwise rules to one another, so the result is $\Gamma$-equivariant.
Different cosets have independent product laws. Thus the product of
the codes preserves probability and has the product of their inverses.
Intersect the countably many good coset domains and their translates.
The resulting domain is invariant and conull. Product continuity
gives topo-finitary regularity; a discrete output reads a finite
set in its one coset.
\end{proof}

\begin{lemma}[Finitary sparse-label homogenization]\label{lem:sparse-code}
Let $H=A_a$ have order $h\geq2$. There is a finite probability base
$(L_H,\lambda_H)$ with $H(\lambda_H)=d\log q/h$ and a finitary
conjugacy over $\Omega_p$
\[
 \Omega_p\times X_H(d)\ \cong\
 \Omega_p\times(L_H,\lambda_H)^G .
\]
\end{lemma}
\begin{proof}
Write $A=H\oplus A'$ and choose distinct $a_j\in A'$, $j\geq1$.
Obtain a uniform priority $\eta_g$ from the biased marker sequence
$(\omega(ga_j))_{j\geq1}$ using the pair extractor from
Lemma~\ref{lem:sparse-start}. For a fixed coset $gH$, the coordinates
$gh_0a_j$, $h_0\in H$, $j\geq1$, are distinct. Hence its $h$
priorities are independent uniforms. Select their unique minimum.
Call the selected set $S_H(\omega)$. It is covariant, has exactly one
element in every $H$-coset, and selects each fixed element with
probability $1/h$. Each selection has a finite marker certificate
because the finite collection of priorities has a strict gap.

Store a coset label only at its selected representative:
\[
 w(g)=
 \begin{cases}
 x(gH),&g\in S_H(\omega),\\
 *,&g\notin S_H(\omega).
 \end{cases}
\]
This is a relative finitary bijection to its constrained image:
recover $x(gH)$ by finding the selected element in the finite coset.
Given the whole marker, the labels at selected elements are
independent uniform $K^d$-labels.

On each line $g\langle s\rangle$ define
\[
 W^g_n=\mathbf1_{\{gs^n\in S_H(\omega)\}},\qquad
 X^g_n=w(gs^n)\quad(n\in\mathbb Z).
\]
The coordinates determining $W^g_n$ belong to $gs^nD$.
These $D$-cosets have different integer heights for different $n$.
Therefore $W^g$ is iid Bernoulli-$1/h$.
Given $W^g$, the variables $X^g_n$ have exactly the sparse product
law in Lemma~\ref{lem:integer-coding}(ii).

Choose a finite base of entropy $c_H=d\log q/h$. Such a base exists
for every finite $c_H$: on a sufficiently large finite alphabet,
entropy varies continuously from zero to the logarithm of its size.
Let $(W,X)\mapsto(W,\varphi_W(X))$ be the relative integer code
from Lemma~\ref{lem:integer-coding}(ii), and set
\[
 z(g)=\varphi_{W^g}(X^g)_0 .
\]
Integer equivariance makes this a consistent rule on each line,
independent of its chosen origin. Left $G$-translation carries
the inputs and lines covariantly to one another, so the rule is
$G$-equivariant.

Here the full conditional law deserves attention. Fix the entire
marker. Every selected position uses a distinct original $H$-coset
label. Consequently the sparse label processes on different lines
are independent under this conditional law, although the lines'
marker processes can be correlated. By disintegration of the
integer isomorphism, for almost every $W$ the fiber code and its
inverse preserve the specified conditional probability laws.
Almost every marker makes every line good, since there are
countably many lines and each $W^g$ has the required iid marginal.
The product of these fiber codes has conditional output law
$\lambda_H^G$, independent of the full marker.

Invert separately on each line, then recover the coset labels.
The joint law and inverse just verified prove a conjugacy over the
entire marker. Each integer output reads finitely many sparse
labels and selection indicators. Each indicator has a finite
marker certificate. A finite union of these inspection sets is
finite. The inverse satisfies the same argument. Countable
intersection of the good domains makes them invariant and conull.
No independence from the exhaustion streams is needed: both
conjugacies fix the marker and preserve each conditional data law.
\end{proof}

\begin{proof}[Proof of Theorem~\ref{thm:all-bases}]
Combine Lemmas~\ref{lem:sparse-start} and~\ref{lem:sparse-code} to obtain
\[
 V_{p,H}^G\cong_{\rm tf}\Omega_p\times\Binf,\qquad
 V_{p,H}=(\{0,1\}\times L_H,\nu_p\times\lambda_H),
\]
where
\[
 H(V_{p,H})=\epsilon_p+\frac{d\log q}{|H|}.
\]
This positive finite entropy can be made arbitrarily small by
taking $p$ small and $H$ large.

The Polish bases of $\Omega_p\times\Binf$ and $\Binf$ both have
infinite entropy. Apply Lemma~\ref{lem:integer-coding}(i) and then
Lemma~\ref{lem:coinduce} with $B=\langle s\rangle$ to obtain
\[
 V_{p,H}^G\cong_{\rm tf}\Binf. \tag{*}
\]
This use of integer coding retains topology; an arbitrary nonatomic
measure-space isomorphism would not suffice.

Now let $(L,\lambda)$ be a nontrivial Polish probability base and
write $b=H(\lambda)\in(0,\infty]$. If $b=\infty$, the same
equal-entropy integer coding and extension immediately give
$L^G\cong_{\rm tf}\Binf$.
If $b<\infty$, choose $V=V_{p,H}$ with entropy $a<b$, and a finite
base $R$ of entropy $b-a$. Equal-entropy integer coding, its
extension to $G$, the conjugacy $(*)$, and one more
infinite-entropy integer coding give
\[
 L^G\cong_{\rm tf}(V\times R)^G
     \cong_{\rm tf}\Binf\times R^G
     \cong_{\rm tf}\Binf .
\]
All good domains can be restricted to invariant conull Borel sets,
and compositions are taken on the preimages of these sets.
Compose the code for one base with the inverse code for the other
to prove the theorem. When both bases are countable and discrete,
continuity of each output label means constancy on a finite
cylinder neighborhood in the conull domain. Thus the final
conjugacy is finitary, even if intermediate alphabets were
nondiscrete.
\end{proof}

\section{Finitely presented groups and a graph formulation}\label{sec:graph}

Lemma~\ref{lem:coinduce} transfers the all-base conclusion to every
countable overgroup of $G$. For the algebraic source in
Appendix~\ref{sec:source}, $N$ is finitely presented. We can preserve a
direct-product construction while making the auxiliary factor finitely
presented. Consider the fixed group
\[
 B=\langle a,u,t\mid
 a^2=1,\ [a,a^t]=1,\ [u,t]=1,\ a^u=aa^t\rangle,
\]
where $x^y=y^{-1}xy$ and $[x,y]=x^{-1}y^{-1}xy$.
This familiar metabelian overgroup of the lamplighter group appears in
Cleary--Riley~\cite{CR}. We give the required facts explicitly.

\begin{lemma}\label{lem:metabelian-stabilizer}
The finitely presented group $B$ is metabelian and contains
$C_2\wr\ZZ$ as the subgroup generated by $a,t$.
\end{lemma}
\begin{proof}
Write $a_j=a^{t^j}$. We first prove $[a_i,a_j]=1$ for all integers
$i,j$. Conjugating by powers of $t$ reduces this to nonnegative
distances $j-i$. Distances zero and one follow from the presentation.
Suppose the claim holds through distance $n\geq1$. Conjugating
$[a_0,a_n]=1$ by $u$ gives
$[a_0a_1,a_na_{n+1}]=1$, since $u$ commutes with $t$.
Every commutator between these four elements except possibly
$[a_0,a_{n+1}]$ is already trivial. Cancelling the other factors
therefore gives $[a_0,a_{n+1}]=1$. Induction proves the claim.
Thus $A_0=\langle a_j:j\in\ZZ\rangle$ is abelian of exponent two.

The relation $a_j^u=a_ja_{j+1}$ shows that conjugation by every
nonnegative power of $u$ carries each $a_j$ into $A_0$.
For any two elements $a^{u^i t^j}$, simultaneously conjugate by a
sufficiently large power of $u$ to make both $u$-exponents
nonnegative. Their images lie in $A_0$ and commute, so the original
pair commutes. The subgroup generated by all these conjugates is
therefore abelian. It is normal under $u,t$ and contains $a$,
so is the normal closure of $a$. The quotient is $\ZZ^2$.
This proves that $B$ is metabelian, and hence amenable.

It remains to exclude extra relations among $a,t$. Over the rational
function field $\mathbb F_2(z)$, the matrices
\[
 a\longmapsto\begin{pmatrix}1&1\\0&1\end{pmatrix},\qquad
 t\longmapsto\begin{pmatrix}z^{-1}&0\\0&1\end{pmatrix},\qquad
 u\longmapsto\begin{pmatrix}(1+z)^{-1}&0\\0&1\end{pmatrix}
\]
satisfy all four defining relations. They send $a_j$ to
$\left(\begin{smallmatrix}1&z^j\\0&1\end{smallmatrix}\right)$.
Distinct Laurent monomials are linearly independent over
$\mathbb F_2$, and the diagonal entry detects each power of $t$.
Thus the homomorphism from the lamplighter presentation into
$\langle a,t\rangle$ is injective as well as onto.
\end{proof}

\begin{corollary}\label{cor:fp}
For the finitely presented algebraic source $N$, the explicit direct
product $N\times B$ has the all-base conclusions of
Theorem~\ref{thm:all-bases}. Here $B$ is a fixed finitely presented
amenable group, independent of $N$. There is also a two-generator
finitely presented group $\Gamma$ with these conclusions.
\end{corollary}
\begin{proof}
The direct product $N\times B$ is finitely presented and contains
$N\times\Lambda$. Apply Lemma~\ref{lem:coinduce}.
Theorem IV of Higman--Neumann--Neumann~\cite{HNN}, including its
finite-relator conclusion, embeds $N\times B$ in a two-generator
finitely presented group $\Gamma$. Apply the same lemma again.
\end{proof}

The group $\Gamma$ is not virtually abelian: it contains the finite
lamp groups $C_2^n$ for all $n$, so its finite subgroups have unbounded
orders. A finitely generated virtually abelian group has a
torsion-free abelian subgroup of finite index, and each finite
subgroup injects into the set of its cosets. Its finite subgroups
therefore have uniformly bounded orders, a contradiction.
Leemann--de la Salle~\cite[Theorem 1.1]{LdS} therefore supplies a
finite-degree Cayley graph $\mathcal G$ of $\Gamma$ whose full
automorphism group is exactly its left translation group.

\begin{corollary}\label{cor:graph}
There is a connected locally finite vertex-transitive graph
$\mathcal G$ on which all nontrivial iid vertex-label processes
with countable discrete bases are finitarily conjugate under
$\operatorname{Aut}(\mathcal G)$. In particular this holds for
every two independent site-percolation processes with parameters
strictly between zero and one.
\end{corollary}
\begin{proof}
Apply Corollary~\ref{cor:fp} on the graphical regular representation
just described. Every finite inspection set lies in a finite
graph ball, so the codes have almost surely finite coding radii.
\end{proof}

The coding can change the geometry of occupied clusters. No
monotonicity, deterministic radius bound, or integrable radius
bound is asserted. For a discrete-valued finitary factor of an iid
process with a countable discrete base, composition with the source
conjugacy gives a finitary realization from Bernoulli-$p$ bits for
every $0<p<1$. The same holds for a topo-finitary discrete-valued
factor of a Polish-base iid process, by composition of continuous maps.
This converts the entire source reversibly before taking the
given factor; it does not supply an inverse to an arbitrary factor.

The distinction is relevant to the factor-of-iid viewpoint in
Lyons~\cite{Lyons}. Mutual Bernoulli factor maps are already known
for every countable nonamenable group by Bowen~\cite{Bowen}.
Here recovery of every input coordinate gives the additional
two-sided conclusion.

\section{What the construction uses}\label{sec:scope}

The auxiliary group $\Lambda=C_2\wr\ZZ$ is fixed independently of the
algebraic input. It is metabelian, hence amenable. The proof uses its
finite lamp subgroups and the height homomorphism; it does not require
a free-product factor. More generally, the same formulas work in any
countable ambient group containing $D=N\times A$ and an infinite-order
element $s$ with $D\cap\langle s\rangle=\{e\}$. This intersection condition
gives both marker-address injectivity and disjointness of the $D$-cosets
along each $s$-line. No step requires the $H_n$ to be normal in the
ambient group.

In the finite-uniform engine, the retained marker contributes exactly
one uniform $K$-coordinate per element of $G$. Its apparent abundance
of randomness inside a $D$-coset
comes from the disjoint marker coordinates indexed by $us^j$.
No independent infinite-base $G$-Bernoulli system is included in the
input.

The remaining input has $d$ coordinates per $H_n$-coset. Its coordinate
density per group element is $d/|H_n|$, which tends to zero. This
decreasing density makes recovery of every coordinate possible even when the inverse
rows have very large finite supports. The choice $k_n=|H_{n+1}|$ offsets
the decreasing density in the extracted Bernoulli factor: each stage
still contributes $\Delta$ ordinary tracks.

The finite groups $H_n$ have orders that may vanish in $K$. We never
average over them. They are used only for finite partitions, orderings,
and coordinate permutations. The only subgroup averages are over the
$J_v$, whose order is invertible by hypothesis.

The core finite-uniform construction uses no entropy classification.
The all-base extension uses precisely the two integer finitary-coding
inputs in Lemma~\ref{lem:integer-coding}. No assertion that zero
Rokhlin entropy alone implies Bernoulli isomorphism is needed.
Continuity is asserted on invariant conull sets. A continuous
conjugacy between the entire two-symbol and four-symbol full shifts
is impossible: it would biject their group-fixed configurations,
of which there are respectively two and four.

\appendix
\section{The precise external algebraic dependency}\label{sec:source}

The algebraic existence input is due to OpenAI. We use \cite{OAI},
\emph{A Counterexample to Kaplansky's Direct-Finiteness Conjecture in
Characteristic Two}, manuscript dated September 23, 2026. We use its
rectangular construction before the subsequent packing into a scalar
one-sided inverse. The implication proved above requires the actual
data of Hypothesis~\ref{hyp:algebra}, not merely an abstract failure of
direct finiteness.

The source's subsection ``Subgroup averages and a rectangular identity''
constructs a countable group $N$ with embedded subgroups, denoted there
by $K_v$, each of order $\kappa=\ell^2$ for an odd prime $\ell$.
Rename them $J_v$ here. Its field $K$ is finite of characteristic two,
so $\kappa$ is invertible. With source vectors $x_v\in K^m$, its matrices
are
\[
 U_{i,v}=-\kappa(x_v)_i p_v,
 \qquad T_{v,i}=p_v(x_v)_i.
\]
The displayed rectangular identities in that subsection are
$UP=U$, $PT=T$, and $UT=I_m$. The source's number $t$ of local subgroups
is our $d$. These formulas also give $C_0\leq\kappa m$ in
\eqref{eq:C}, although our proof needs only finiteness of $C_0$.

The strict inequality needed here is slightly stronger than the weak
inequality displayed in the source's incidence criterion. It follows
from its terminating prescription. Write $Q_0$ for the number of affine
points in that prescription. The selected parameters satisfy
\[
 m\geq\frac{2Q_0}{(\ell+1)^2},\qquad d=Q_0-m.
\]
Consequently
\[
 m(\ell^2+1)
 \geq \frac{2(\ell^2+1)}{(\ell+1)^2}Q_0
 >Q_0,
\]
because
$2(\ell^2+1)-(\ell+1)^2=(\ell-1)^2>0$.
Rearranging gives $d<\ell^2m=\kappa m$, as required.

Thus the cited existence construction supplies all of
Hypothesis~\ref{hyp:algebra}, and Theorem~\ref{thm:main} applies to its
data. The existence of that incidence construction and the embedding
of its local groups are the external results invoked here. This note
does not reproduce their proofs or independently certify a formal
verification of them. Every subsequent operator, probability,
exhaustion, and absorption argument is given in the body of this note.

The same source's local-group embedding construction gives a finite
presentation of this $N$: it has finitely many generators and the finitely
many local multiplication-table relations of length at most three.
Corollary~\ref{cor:fp} uses this finite presentation, together with the
classical embedding theorems cited there.

\begin{remark}[The algebraic input can be binary]\label{rem:binary}
For a characteristic-two instance of Hypothesis~\ref{hyp:algebra},
let $f=[K:\mathbb F_2]$. Choose an $\mathbb F_2$-basis of $K$ and
represent multiplication by each field element as an $f$-by-$f$
binary matrix. Applying this representation to every coefficient
gives a unital algebra homomorphism
$K[N]\to\Mat_f(\mathbb F_2[N])$. Expanding $U,T,P$ into blocks
therefore preserves their three identities. Each average $p_v$
has coefficients in $\mathbb F_2$ because $\kappa$ is odd, so its
block is $p_v I_f$. The expanded diagonal averages correspond to
repeating each $J_v$ exactly $f$ times. The new dimensions are
$d'=fd$, $m'=fm$, and $d'<\kappa m'$ still holds.
Thus the construction may be carried out with binary data.
\end{remark}

\begin{thebibliography}{99}
\bibitem{OAI}
OpenAI,
\emph{A Counterexample to Kaplansky's Direct-Finiteness Conjecture in
Characteristic Two}, manuscript dated September 23, 2026.
\href{https://github.com/openai/math/tree/adc7f1241b42e322a6451854ab7e4b4c146bf78a/preprints/A-Counterexample-to-Kaplanskys-Direct-Finiteness-Conjecture-in-Characteristic-Two-September-23-2026}{Manuscript and source files}.
The source revision is \texttt{adc7f1241b42e322a6451854ab7e4b4c146bf78a}.
The specific source passages used here are
\href{https://raw.githubusercontent.com/openai/math/adc7f1241b42e322a6451854ab7e4b4c146bf78a/preprints/A-Counterexample-to-Kaplanskys-Direct-Finiteness-Conjecture-in-Characteristic-Two-September-23-2026/build/sections/02-criterion.tex}{\texttt{02-criterion.tex}},
subsection ``Subgroup averages and a rectangular identity'', and
\href{https://raw.githubusercontent.com/openai/math/adc7f1241b42e322a6451854ab7e4b4c146bf78a/preprints/A-Counterexample-to-Kaplanskys-Direct-Finiteness-Conjecture-in-Characteristic-Two-September-23-2026/build/sections/05-prescription.tex}{\texttt{05-prescription.tex}},
subsection ``Finite search for the incidence data''; and
\href{https://raw.githubusercontent.com/openai/math/adc7f1241b42e322a6451854ab7e4b4c146bf78a/preprints/A-Counterexample-to-Kaplanskys-Direct-Finiteness-Conjecture-in-Characteristic-Two-September-23-2026/build/sections/02a-embedding.tex}{\texttt{02a-embedding.tex}},
for the local subgroup embeddings and finite presentation.
Accessed October 25, 2026.
\bibitem{MS}
T. Meyerovitch and Y. Spinka,
\emph{Entropy-efficient finitary codings},
\href{https://arxiv.org/abs/2201.06542}{arXiv:2201.06542v1},
Theorems 3.2 and 3.15; version consulted October 25, 2026.
\bibitem{CR}
S. Cleary and T. R. Riley,
\emph{A finitely presented group with unbounded dead-end depth},
Proc. Amer. Math. Soc. \textbf{134} (2006), 343--349.
\href{https://pi.math.cornell.edu/~riley/papers/Dead-End_Depth/deadend_corrected.pdf}
{Corrected author version, revised April 5, 2007}, Section 1.
\bibitem{HNN}
G. Higman, B. H. Neumann, and H. Neumann,
\emph{Embedding theorems for groups},
J. London Math. Soc. \textbf{24} (1949), 247--254, Theorem IV.
\href{https://doi.org/10.1112/jlms/s1-24.4.247}{doi:10.1112/jlms/s1-24.4.247}.
\bibitem{LdS}
P.-H. Leemann and M. de la Salle,
\emph{Cayley graphs with few automorphisms: the case of infinite groups},
Ann. H. Lebesgue \textbf{5} (2022), 73--92, Theorem 1.1.
\href{https://doi.org/10.5802/ahl.118}{doi:10.5802/ahl.118}.
\bibitem{Lyons}
R. Lyons, \emph{Factors of IID on trees},
Combin. Probab. Comput. \textbf{26} (2017), 285--300.
\href{https://rdlyons.pages.iu.edu/pdf/fiid.pdf}{Author's version};
\href{https://rdlyons.pages.iu.edu/errata/fiid.pdf}{2020 erratum}.
\bibitem{Bowen}
L. Bowen, \emph{Finitary random interlacements and the Gaboriau--Lyons problem},
\href{https://arxiv.org/abs/1707.09573}{arXiv:1707.09573}.
\bibitem{Seward}
B. Seward,
\emph{Bernoulli shifts with bases of equal entropy are isomorphic},
\href{https://arxiv.org/abs/1805.08279}{arXiv:1805.08279}, Theorem 1.1.
\bibitem{BowenSofic}
L. Bowen,
\emph{Measure conjugacy invariants for actions of countable sofic groups},
J. Amer. Math. Soc. \textbf{23} (2010), 217--245.
\href{https://arxiv.org/abs/0804.3582}{arXiv:0804.3582}, Theorem 1.1.
\bibitem{Astra}
Astra research archive, public research notes in
\emph{astra-research-knowledge}:
\href{https://github.com/harjassand/astra-research-knowledge/blob/39cdea532f19217f633f70315d92c314a6ce5e3e/cards/N252-finite-defect-Bernoulli-extraction.txt}{N252, finite defect extraction};
\href{https://github.com/harjassand/astra-research-knowledge/blob/0146b0d4d6ea02edf757bdb69a517378bf801578/cards/N279-finite-field-Bernoulli-premise-repair.txt}{N279, finite-field premise repair}; and
\href{https://github.com/harjassand/astra-research-knowledge/blob/142ac61f9b738eeddd466e9cc591b607689a542d/web/pages/d-377e83d87f921de6.html}{N270, random basis and collapse criterion}.
Versions consulted October 25, 2026. These are public partial or conditional
research developments, cited for comparison rather than as inputs to the proof.
\bibitem{Sauers}
SauersML, public research notes in \emph{group-approximation}:
\href{https://github.com/SauersML/group-approximation/blob/0d15e308d317eeed73fbdb4e006797a4d6b080d7/research/bernoulli-isomorphism-from-approximate-right-inverses.md}{Bernoulli isomorphism from approximate right inverses}; and
\href{https://github.com/SauersML/group-approximation/blob/ba6d9b822153b7feb2de9ce2caead1c30a18050f/research/zero-supremum-endomorphisms-have-approximate-right-inverses.md}{the status of the approximate-right-inverse prerequisite}.
Versions consulted October 25, 2026. The conjugacy criterion retains an
additional hypothesis; it is cited for comparison, not used as a proof input.
\end{thebibliography}
\end{document}
